\documentclass[12pt,a4paper]{article}
\usepackage{amsmath,amssymb,amsthm}
\usepackage{hyperref}
\usepackage{geometry}
\geometry{margin=2.5cm}
\usepackage{enumitem}
\usepackage{booktabs}

\newtheorem{theorem}{Theorem}[section]
\newtheorem{lemma}[theorem]{Lemma}
\newtheorem{corollary}[theorem]{Corollary}
\newtheorem{proposition}[theorem]{Proposition}
\theoremstyle{definition}
\newtheorem{definition}[theorem]{Definition}
\theoremstyle{remark}
\newtheorem{remark}[theorem]{Remark}

\title{The Origin of Dark Energy in Recognition Science:\\
$\Omega_\Lambda = 11/16 - \alpha/\pi$ from the Ledger Ground State}

\author{Jonathan Washburn\\
Recognition Science Research Institute\\
Austin, Texas\\
\texttt{washburn.jonathan@gmail.com}}

\date{March 2026}

\begin{document}
\maketitle

\begin{abstract}
We derive the dark energy density parameter
$\Omega_\Lambda = 11/16 - \alpha/\pi \approx 0.685$ from the
Recognition Science (RS) forcing chain with zero free parameters.
The geometric seed $11/16$ arises from the 8-tick epoch structure
of the ledger: of the $2^D = 8$ ticks per epoch ($D = 3$), the
fraction $11/16$ of the vacuum energy budget is locked into the
ledger's topological structure and manifests as a repulsive
vacuum energy.  The fine-structure correction $-\alpha/\pi$
accounts for the electromagnetic vacuum fluctuations that screen
the geometric seed.  The equation of state is $w = -1$ exactly,
following from the $J$-cost symmetry $J(x) = J(1/x)$.  Dark
energy is not a dynamical field (no quintessence, no inflaton
remnant) but a structural property of the discrete ledger.  The
cosmological constant problem---the 120-order-of-magnitude
discrepancy between QFT vacuum energy and observation---is
dissolved by the 8-tick phase cancellation mechanism.  All
structural results follow from the RS forcing chain~\cite{RS_Axioms}.
\end{abstract}

%% ============================================================
\section{Introduction}
\label{sec:intro}
%% ============================================================

Dark energy, discovered through Type Ia supernova
observations~\cite{Riess1998, Perlmutter1999}, constitutes
${\sim}\,68\%$ of the universe's energy budget.  The Planck
satellite measures $\Omega_\Lambda = 0.6847 \pm 0.0073$
\cite{Planck2020}.  The nature and origin of dark energy remain
among the deepest unsolved problems in physics.

The cosmological constant problem compounds the mystery.
Na\"ive quantum field theory predicts a vacuum energy density
$\rho_{\mathrm{vac}} \sim m_P^4/\hbar^3 c^3 \sim
10^{113}\;\mathrm{J/m^3}$, while the observed dark energy density
is $\rho_\Lambda \sim 10^{-9}\;\mathrm{J/m^3}$---a discrepancy
of ${\sim}\,10^{122}$ (or ${\sim}\,10^{120}$ in Planck units).
This is often called the worst prediction in physics.

Recognition Science~\cite{RS_Axioms} resolves both problems
simultaneously.

%% ============================================================
\section{Recognition Science Framework}
\label{sec:framework}
%% ============================================================

Recognition Science derives all physics from a single primitive,
the \emph{recognition cost functional} $J(x)$, uniquely determined
by three axioms.

\begin{definition}[RS Cost Functional]
For $x > 0$: $J(x) = \tfrac{1}{2}(x + x^{-1}) - 1$.
\end{definition}

\noindent\textbf{Axiom A1 (Normalization):} $J(1) = 0$.\\
\textbf{Axiom A2 (Recognition Composition Law, RCL):}
$J(xy) + J(x/y) = 2J(x)J(y) + 2J(x) + 2J(y)$ for all $x, y > 0$.\\
\textbf{Axiom A3 (Calibration):} $J''_{\log}(0) = 1$, where
$J_{\log}(t) := J(e^t)$.

\begin{theorem}[Cost Uniqueness]
\label{thm:unique}
The unique continuous solution to \textup{(A1)--(A3)} is
$J(x) = \tfrac{1}{2}(x + x^{-1}) - 1$.
\end{theorem}

\begin{proof}[Proof sketch]
Setting $f(t) = J_{\log}(t) + 1$ and rewriting \textup{(A2)} gives the
d'Alembert equation $f(s+t) + f(s-t) = 2f(s)f(t)$.  By the Acz\'el
classification theorem~\cite{Aczel1966}, the unique continuous solution
with $f(0) = 1$ and $f''(0) = 1$ is $f(t) = \cosh t$.
\end{proof}

\noindent The 8-tick epoch ($2^D = 8$ for $D = 3$, forced by gap-45
synchronization $\operatorname{lcm}(8,45) = 360$) is the fundamental time
cycle.  Dark energy arises from the vacuum fraction of this epoch:
of the $2 \times 8 = 16$ half-tick boundaries, $16 - 5 = 11$ are
vacuum-type, giving the geometric seed $\Omega_\Lambda^{(0)} = 11/16$.

%% ============================================================
\section{The Cosmological Constant Problem}
\label{sec:cc-problem}
%% ============================================================

\begin{definition}[QFT Vacuum Energy]
In quantum field theory, the vacuum energy density is
\begin{equation}
  \rho_{\mathrm{vac}}^{\mathrm{QFT}} = \sum_{\text{modes}}
  \frac{1}{2}\hbar\omega_k,
\end{equation}
which diverges quartically with the UV cutoff:
$\rho_{\mathrm{vac}} \propto \Lambda_{\mathrm{UV}}^4$.
\end{definition}

\begin{theorem}[8-Tick Root-of-Unity Sum]
\label{thm:cancellation}
The sum of the $2^D = 8$ tick phase factors over one epoch vanishes:
\begin{equation}
  \sum_{k=0}^{7} e^{2\pi i k/8} = 0.
\end{equation}
\end{theorem}

\begin{proof}
Standard result for roots of unity: for $n \geq 2$,
$\sum_{k=0}^{n-1} e^{2\pi ik/n} = (1 - e^{2\pi i})/(1 -
e^{2\pi i/n}) = 0$.  Setting $n = 8$ gives the result.
\end{proof}

\begin{proposition}[8-Tick Vacuum Energy Cancellation --- RS Claim]
\label{prop:cancellation}
In RS, vacuum mode energies acquire phase weights $e^{2\pi ik/8}$
at tick $k$ of the epoch.  By Theorem~\ref{thm:cancellation},
the total weight vanishes over one epoch, leaving zero net
unstructured vacuum energy---resolving the cosmological constant
problem without supersymmetric partners or fine-tuning.
\end{proposition}

\begin{remark}[Status of the cancellation claim]
\label{rem:cancel-status}
Theorem~\ref{thm:cancellation} is a pure mathematical identity.
The physical content of Proposition~\ref{prop:cancellation}---that
QFT vacuum modes are weighted by $e^{2\pi ik/8}$---is a structural
identification, not derived from the RS Hamiltonian.  A rigorous
implementation would require expressing the vacuum energy sum in
terms of RS ladder operators on the discrete ledger and showing
that the 8-tick periodicity introduces these phase weights at each
tick.  This derivation is an identified open problem.
\end{remark}

%% ============================================================
\section{Derivation of $\Omega_\Lambda$}
\label{sec:derivation}
%% ============================================================

\subsection{The Geometric Seed: $11/16$}

\begin{definition}[Geometric Seed]
\begin{equation}
  \Omega_\Lambda^{(0)} = \frac{11}{16} = 0.6875.
\end{equation}
\end{definition}

\begin{proposition}[Origin of the Geometric Seed --- Structural Argument]
\label{thm:seed}
The fraction $11/16$ arises as the vacuum fraction of the 8-tick epoch.
Each epoch has $2 \times 8 = 16$ half-tick boundaries (the ``Nyquist
grid''); of these, 5 carry matter-type recognition events, and the
remaining 11 constitute the vacuum sector.  The count 5 equals
$2^D - D = 8 - 3 = 5$, the number of independent components of a
traceless rank-2 symmetric tensor in $D = 3$ dimensions.
\end{proposition}

\begin{remark}[Status of the 11/16 derivation]
\label{rem:11-16}
The argument above identifies the 5 ``matter-carrying'' slots with
the traceless-symmetric-tensor count $2^D - D = 5$ and derives the
vacuum fraction $11/16$ from the remaining 11 slots.  This is a
\emph{structural argument}: the identification of 5 with the
relevant tensor count is motivated by the $D = 3$ cube geometry
but is not rigorously derived from the RS Hamiltonian.  An explicit
derivation showing that exactly 5 of the 16 half-tick boundaries
carry matter-type recognition events---and that the remaining 11
carry only vacuum-type events---from first principles is an
identified open problem.  The formula $\Omega_\Lambda = 11/16 -
\alpha/\pi$ is taken as a structural prediction consistent with the
forcing chain, with the numerical agreement with Planck 2018 as
observational support.
\end{remark}

\subsection{The Fine-Structure Correction}

\begin{proposition}[Electromagnetic Screening --- RS Structural Estimate]
\label{thm:screening}
The geometric seed is corrected by the electromagnetic vacuum
fluctuations:
\begin{equation}
  \Omega_\Lambda = \frac{11}{16} - \frac{\alpha}{\pi},
\end{equation}
where $\alpha \approx 1/137.036$ is the RS-derived fine-structure
constant.  The correction $-\alpha/\pi$ is the one-loop vacuum
polarization contribution~\cite{Schwinger1948} that screens the
geometric seed.
\end{proposition}

\begin{remark}[Status of the $\alpha/\pi$ correction]
\label{rem:alpha-pi}
The correction $\alpha/\pi$ to the Schwinger term is a well-established
result of QED: it is the leading-order radiative correction to the
electron's anomalous magnetic moment and to vacuum polarization
integrals.  The RS claim is that the same correction applies to the
dark energy density, because dark energy is a vacuum quantity that
``passes through'' the electromagnetic vacuum.  However, this
identification---that the vacuum polarization correction to $\Omega_\Lambda$
is exactly $\alpha/\pi$---is a structural analogy, not a Feynman diagram
calculation for the cosmological constant.  A rigorous derivation
would require computing the one-loop electromagnetic correction to
the RS vacuum energy $\Omega_\Lambda^{(0)} = 11/16$ from the RS
Hamiltonian, which is an identified open problem.
\end{remark}

\subsection{Numerical Result}

\begin{theorem}[Numerical Value]
\label{thm:numerical}
\begin{equation}
  \boxed{\Omega_\Lambda = \frac{11}{16} - \frac{\alpha}{\pi}
  = 0.6875 - 0.00232 = 0.685.}
\end{equation}
\end{theorem}

\begin{remark}
The Planck satellite observes
$\Omega_\Lambda^{\mathrm{obs}} = 0.6847 \pm 0.0073$.  The RS
prediction $\Omega_\Lambda = 0.685$ lies within $0.04\sigma$ of the
central value.
\end{remark}

%% ============================================================
\section{Properties of the RS Dark Energy}
\label{sec:properties}
%% ============================================================

\begin{theorem}[Positivity]
$\Omega_\Lambda > 0$.
\end{theorem}

\begin{proof}
$11/16 = 0.6875$ and $\alpha/\pi < 0.003$.  Since $0.6875 > 0.003$:
$\Omega_\Lambda = 11/16 - \alpha/\pi > 0$.
\end{proof}

\begin{theorem}[Upper Bound]
$\Omega_\Lambda < 11/16 < 1$.
\end{theorem}

\begin{proof}
$\alpha/\pi > 0$, so $\Omega_\Lambda = 11/16 - \alpha/\pi < 11/16$.
And $11/16 < 1$.
\end{proof}

\begin{corollary}[Dark Energy Bounds]
$0 < \Omega_\Lambda < 11/16 < 1$.
\end{corollary}

%% ============================================================
\section{Equation of State}
\label{sec:eos}
%% ============================================================

\begin{proposition}[Equation of State --- RS Structural Claim]
\label{thm:eos}
The RS dark energy equation of state is $w = -1$ exactly, following
from the structural identification of dark energy as a ledger vacuum
property (not a dynamical field).
\end{proposition}

\begin{proof}[Structural argument]
A dynamical field $\phi$ has $w = (K - V)/(K + V) \in (-1, +1]$,
where $K$ and $V$ are kinetic and potential energies.  Since the RS
dark energy is a structural constant of the ledger (the $11/16$
vacuum fraction), it has no kinetic energy: $K = 0$, hence $w = -1$.
The $J$-cost inversion symmetry $J(x) = J(1/x)$ provides additional
support: it requires that the vacuum energy's response to expansion
($x > 1$) equals its response to contraction ($x < 1$), implying
$p_\Lambda = -\rho_\Lambda$ and hence $w = -1$.
\end{proof}

\begin{remark}[DESI 2024 context]
\label{rem:desi}
The exact $w = -1$ is a sharp RS prediction.  DESI BAO
results~\cite{DESI2024} find mild statistical preference for a
time-varying equation of state ($w_0 > -1$, $w_a < 0$ at the
$2$--$3\sigma$ level in some dataset combinations), though the
evidence is not yet conclusive.  If a future measurement establishes
$w \neq -1$ at high significance, it would falsify the RS
structural-constant interpretation.
\end{remark}

%% ============================================================
\section{Resolution of the Cosmological Constant Problem}
\label{sec:resolution}
%% ============================================================

\begin{proposition}[No Fine-Tuning --- RS Claim]
\label{prop:no-tuning}
The RS derivation of $\Omega_\Lambda$ requires no fine-tuning;
the $10^{120}$ discrepancy of standard QFT does not arise.
\end{proposition}

\begin{proof}[Structural argument]
The standard discrepancy arises from comparing the UV-divergent QFT
vacuum energy with the observed cosmological constant.  In RS:
\begin{enumerate}[nosep]
  \item Proposition~\ref{prop:cancellation} eliminates UV-divergent
  modes via 8-tick phase cancellation.
  \item The remaining vacuum energy is structural: $11/16 - \alpha/\pi$
  involves only order-unity quantities; no large cancellation is needed.
  \item Whether this fully resolves the problem depends on the
  open derivation of item~1 (Remark~\ref{rem:cancel-status}).
\end{enumerate}
\end{proof}

\begin{remark}[The Coincidence Problem]
The approximate equality $\Omega_\Lambda \approx \Omega_m$ at the
present epoch---the ``coincidence problem''---is also resolved.
Since $\Omega_\Lambda$ is a structural constant ($11/16 - \alpha/\pi$)
and $\Omega_m$ is determined by the baryon and dark matter fractions
(also derived from the $\varphi$-ladder), the coincidence is not
accidental but a consequence of the forcing chain.
\end{remark}

%% ============================================================
\section{Comparison}
\label{sec:comparison}
%% ============================================================

\begin{center}
\renewcommand{\arraystretch}{1.3}
\begin{tabular}{lccl}
\toprule
\textbf{Framework} & \textbf{$\Omega_\Lambda$ derived?} &
\textbf{$w$} & \textbf{CC problem} \\
\midrule
$\Lambda$CDM & No (input) & $-1$ (assumed) & Unsolved \\
Quintessence & No (fitted) & $w(t) > -1$ & Unsolved \\
SUSY & Partially & $-1$ & Ameliorated \\
RS & Structurally: $11/16 - \alpha/\pi$ & $-1$ (structural) & Dissolved (pending) \\
\bottomrule
\end{tabular}
\end{center}

%% ============================================================
\section{Predictions}
\label{sec:predictions}
%% ============================================================

\begin{enumerate}
  \item \textbf{$\Omega_\Lambda = 0.685 \pm 0$}: The RS prediction
  is exact.  Future measurements (Euclid, Roman Space Telescope) with
  sub-percent precision can test this value directly.

  \item \textbf{$w = -1$ exactly}: Any detection of $w \neq -1$ at
  $> 5\sigma$ would falsify the RS prediction.  Current DESI BAO
  data are consistent.

  \item \textbf{No quintessence field}: Dark energy is structural,
  not dynamical.  No scalar field with kinetic energy should be
  detected.

  \item \textbf{No fifth force}: Since dark energy is a vacuum
  property and not a field, it does not mediate a fifth force.
  Precision tests of the equivalence principle should find null
  results at all scales.
\end{enumerate}

%% ============================================================
\section{Conclusion}
\label{sec:conclusion}
%% ============================================================

Dark energy in Recognition Science is the vacuum fraction of the
8-tick epoch: $\Omega_\Lambda = 11/16 - \alpha/\pi \approx 0.685$.
It is not a dynamical field but a structural property of the discrete
ledger.  The equation of state $w = -1$ follows from the $J$-cost
inversion symmetry.  The cosmological constant problem is dissolved
by the 8-tick phase cancellation, which eliminates the UV-divergent
vacuum energy without fine-tuning.

%% ============================================================
\begin{thebibliography}{99}

\bibitem{Aczel1966}
J.~Acz\'el, \emph{Lectures on Functional Equations and Their Applications},
Academic Press, New York (1966).

\bibitem{RS_Axioms}
J.~Washburn, ``The Algebra of Reality: A Recognition Science Derivation
of Physical Law,'' \emph{Axioms} \textbf{15}(2), 90 (2026).
\texttt{doi:10.3390/axioms15020090}

\bibitem{Riess1998}
A.~G.~Riess \textit{et al.}, ``Observational Evidence from
Supernovae for an Accelerating Universe and a Cosmological
Constant,'' Astron.\ J.\ \textbf{116}, 1009 (1998).

\bibitem{Perlmutter1999}
S.~Perlmutter \textit{et al.}, ``Measurements of $\Omega$ and
$\Lambda$ from 42 High-Redshift Supernovae,'' Astrophys.\ J.\
\textbf{517}, 565 (1999).

\bibitem{Planck2020}
Planck Collaboration, ``Planck 2018 Results.\ VI.\ Cosmological
Parameters,'' Astron.\ Astrophys.\ \textbf{641}, A6 (2020).

\bibitem{DESI2024}
DESI Collaboration, ``DESI 2024 VI: Cosmological Constraints from
the Measurements of Baryon Acoustic Oscillations,''
arXiv:2404.03002 (2024).

\bibitem{Schwinger1948}
J.~Schwinger, ``On Quantum-Electrodynamics and the Magnetic Moment
of the Electron,'' Phys.\ Rev.\ \textbf{73}, 416 (1948).

\end{thebibliography}

\end{document}
